\section*{الفصل \ref{ch:g6:water-underground} --- الماء تحت الأرض}

\begin{solution}{exo:g6:water-underground:1}
\omterm{def:g6:water-underground:permeability}{النفوذة}: \omterm{def:g3:stones-pebbles-sand:sand}{الرمل}، \omterm{def:g3:stones-pebbles-sand:pebble}{والحصباء}، والطباشير. \omterm{def:g6:water-underground:permeability}{الكتيمة}: الطين، والغرانيت الخالي من الشقوق،
والسجّيل.
\end{solution}

\begin{solution}{exo:g6:water-underground:2}
(أ) \omterm{def:g6:water-underground:water-table}{منسوب المياه الجوفية}؛ (ب) \omterm{def:g6:water-underground:aquifer}{خزان جوفي}؛ (ج) \omterm{def:g6:water-underground:spring}{ينبوع}؛ (د) \omterm{def:g6:water-underground:spring}{بئر}.
\end{solution}

\begin{solution}{exo:g6:water-underground:3}
(أ) يتبخر عائدًا إلى الهواء؛ ويجري على السطح إلى
الجداول؛ ويتسرب في الأرض (\omterm{def:g6:water-underground:infiltration}{التسرب}). (ب) في موقف السيارات: فسطحه
الصلب يكاد لا يسمح للماء \omterm{def:g6:water-underground:infiltration}{بالتسرب}، فيجري الماء على السطح. أما في
المرج، فالتربة والعشب يسمحان لقسم كبير منه \omterm{def:g6:water-underground:infiltration}{بالتسرب}.
\end{solution}

\begin{solution}{exo:g6:water-underground:4}
(أ) خطأ: الطين \omterm{def:g6:water-underground:porosity}{مسامي} جدًا لكنه \omterm{def:g6:water-underground:permeability}{كتيم}، لأن مسامه دقيقة جدًا.
(ب) صواب. (ج) خطأ: فهو يتبع شكل الأرض، بلطف أكبر، ويرتفع
وينخفض مع الفصول. (د) صواب: نحو 30~\% من كل \omterm{def:g2:water-on-earth:salt-fresh}{الماء العذب}
\omterm{def:g6:water-underground:groundwater}{مياه جوفية}، مقابل نحو ربع واحد في المئة في \omterm{def:g1:landscapes-around-us:water}{البحيرات}.
\end{solution}

\begin{solution}{exo:g6:water-underground:5}
$10 \times 25 \div 100 = 2.5$: تستطيع أن تصب 2.5~لتر قبل أن يفيض.
$10 \times 22 \div 100 = 2.2$: سيُصرَّف 2.2~لتر عبر الثقب.
\end{solution}

\begin{solution}{exo:g6:water-underground:6}
(أ) $2 \times 1 \times 1 = 2$~متر مكعب، أي 2000~لتر.
(ب) $2000 \times 20 \div 100 = 400$~لتر. (ج) $2000 \times 18 \div 100 =
360$~لترًا تُصرَّف؛ و$400 - 360 = 40$~لترًا تبقى عالقة.
\end{solution}

\begin{solution}{exo:g6:water-underground:7}
(أ) \omterm{def:g6:water-underground:spring}{الينبوع} على \omterm{def:g4:world-physical-map:altitude}{ارتفاع} نحو 203~أمتار، \omterm{def:g1:landscapes-around-us:water}{والنهر} على \omterm{def:g4:world-physical-map:altitude}{ارتفاع} نحو 132~مترًا. (ب) \omterm{def:g6:water-underground:spring}{الينبوع}
أعلى من \omterm{def:g1:landscapes-around-us:water}{النهر} بنحو 70~مترًا: فماؤه يجري نزولًا على المنحدر إلى \omterm{def:g1:landscapes-around-us:water}{النهر}.
(ج) داخل السفح الطباشيري يكون \omterm{def:g6:water-underground:water-table}{منسوب المياه الجوفية} تحت سطح الأرض؛ ولا يلتقي
بالسطح إلا حيث تبرز على المنحدر قاعدة الطباشير، المستندة إلى السجّيل
\omterm{def:g6:water-underground:permeability}{الكتيم}.
\end{solution}

\begin{solution}{exo:g6:water-underground:8}
(أ) بعد موسم الأمطار: $401 - 336 = 65$~مترًا؛ وبعد موسم \omterm{def:g5:natural-hazards:drought}{الجفاف}:
$401 - 314 = 87$~مترًا. (ب) $336 - 314 = 22$~مترًا.
\end{solution}

\begin{solution}{exo:g6:water-underground:9}
أعمق. \omterm{def:g6:water-underground:water-table}{فمنسوب المياه الجوفية} يتبع شكل الأرض، لكن بلطف أكبر: ففي
قمة التل يكون أعلى منه تحت \omterm{def:g6:water-underground:spring}{البئر}~A، لكن الأرض ترتفع
أكثر بكثير، فيقع الماء على عمق أكبر تحت السطح (في المقطع،
على عمق يزيد على 110~أمتار).
\end{solution}

\begin{solution}{exo:g6:water-underground:10}
\omterm{def:g6:water-underground:porosity}{مسام} الطين دقيقة جدًا: فيبقى الماء عالقًا بالرقائق ولا يستطيع أن
يجري عبر الطين نحو \omterm{def:g6:water-underground:spring}{البئر}. فالطين \omterm{def:g6:water-underground:porosity}{مسامي} لكنه \omterm{def:g6:water-underground:permeability}{كتيم}، لذلك
لا يُصرَّف منه إلا قليل من الماء (20~لترًا من أصل 500).
\end{solution}

\begin{solution}{exo:g6:water-underground:11}
(أ) $30.1 \div 0.26 \approx 116$: أكثر بما يزيد على مئة مرة.
(ب) $100 - 69 = 31$: يبقى نحو 31~\%، \omterm{def:g6:water-underground:groundwater}{والمياه الجوفية} وحدها نحو 30
منها. (ج) إنها أكبر مخزون \omterm{def:g2:water-on-earth:salt-fresh}{للماء العذب} السائل بفارق كبير؛ أما الجليد
فمتجمد وبعيد في معظمه، في المناطق القطبية وعلى \omterm{def:g1:landscapes-around-us:landform}{الجبال}. ثم إن
\omterm{def:g6:water-underground:groundwater}{المياه الجوفية} توجد في كل مكان تقريبًا تحت أقدامنا.
\end{solution}

\begin{solution}{exo:g6:water-underground:12}
البئران تستمدان الماء من \omterm{def:g6:water-underground:aquifer}{الخزان الجوفي} الطباشيري نفسه. فقد أخرجت المضخة القوية الماء
أسرع مما يعيد \omterm{def:g4:clouds-rain-snow:precipitation}{المطر} ملء الطباشير، فانخفض \omterm{def:g6:water-underground:water-table}{منسوب المياه الجوفية} في كل
ما حولها، حتى صار تحت قاع \omterm{def:g6:water-underground:spring}{البئر} الثانية، فجفت. ويستطيع المزارعان
أن يتفقا على ضخ كمية أقل، ولا سيما في الصيف، أو أن يتقاسما بئرًا واحدة، أو أن يعمّقا
\omterm{def:g6:water-underground:spring}{البئر} الثانية؛ والأهم أن لا يأخذا من الماء أكثر مما يعيده
\omterm{def:g4:clouds-rain-snow:precipitation}{المطر}.
\end{solution}

\begin{solution}{pb:g6:water-underground:1}
\textbf{1.} تقوم المزرعة على الطباشير؛ وتحته السجّيل. والطباشير
\omterm{def:g6:water-underground:permeability}{نفوذ}، أما السجّيل فلا.

\textbf{2.} يتسرب \omterm{def:g4:clouds-rain-snow:precipitation}{المطر} في الطباشير ويتوقف عند السجّيل. وحيث تبرز
قاعدة الطباشير على المنحدر، يلتقي \omterm{def:g6:water-underground:water-table}{منسوب المياه الجوفية} بسطح الأرض
فيخرج الماء: هذا هو \omterm{def:g6:water-underground:spring}{الينبوع}.

\textbf{3.} $401 - 336 = 65$~مترًا تحت سطح الأرض.

\textbf{4.} قاع \omterm{def:g6:water-underground:spring}{البئر} على ارتفاع $401 - 100 = 301$~متر، تحت
\omterm{def:g6:water-underground:water-table}{منسوب المياه الجوفية} الواقع على \omterm{def:g4:world-physical-map:altitude}{ارتفاع} 314~مترًا: فلا تزال \omterm{def:g6:water-underground:spring}{البئر} تصل إلى الماء، بفارق $314 - 301 =
13$~مترًا.

\textbf{5.} $100 \times 50 \times 2 = 10\,000$~متر مكعب.

\textbf{6.} $10\,000 \times 20 \div 100 = 2000$~متر مكعب.

\textbf{7.} $10\,000 \times 18 \div 100 = 1800$~متر مكعب تُصرَّف؛
و$2000 - 1800 = 200$~متر مكعب تبقى عالقة.

\textbf{8.} $1800 \times 1000 = 1\,800\,000$~لتر.

\textbf{9.} $18 \times 1000 = 18\,000$~لتر في اليوم.

\textbf{10.} $18 \times 7 = 126$~مترًا مكعبًا في الأسبوع.

\textbf{11.} \omterm{def:g4:clouds-rain-snow:precipitation}{مطر} الموسم الرطب يتسرب في الطباشير ويعيد ملأه
كل سنة (ما دامت المزارع لا تضخ أكثر مما يعيده \omterm{def:g4:clouds-rain-snow:precipitation}{المطر}).

\textbf{12.} $1800 \div 18 = $~\textbf{100 يوم}: أكثر قليلًا من ثلاثة
أشهر دون أي \omterm{def:g4:clouds-rain-snow:precipitation}{مطر}.
\end{solution}
